% ============================================================
% Thesis Chapter: Learning Controllers for Energy-Aware UPF Orchestration
% ============================================================
% Compile with:  pdflatex chapter_controllers.tex  (run twice for refs)
% Figures live in  figures/ .
% Diagrams marked [DRAW.IO PLACEHOLDER] are to be drawn by hand.
%
% NUMBERS: all quantitative claims use the post-fix (UPF_NDT v0.4.0)
% evaluation. Per-controller means are from
%   reports/phase-7/multiseed_summary_v04twin.json
% and the comparison note reports/phase-7/switching_fix_reeval.md.
% All figures regenerated under the v0.4 twin in the paper font (Latin Modern,
% usetex) by research/phase7/generate_chapter_figures.py (4 data figures) and
% generate_chapter_heatmaps.py (2 rollout heatmaps). Paired bootstrap recomputed
% under v0.4: reports/phase-7/paired_bootstrap_mappo_vs_ippo_v04.json.
% ============================================================

\documentclass[12pt, a4paper]{report}

% -- Packages ------------------------------------------------------------
\usepackage[T1]{fontenc}
\usepackage[utf8]{inputenc}
\usepackage{lmodern}
\usepackage{amsmath, amssymb}
\usepackage{booktabs}
\usepackage{graphicx}
\usepackage{subcaption}
\usepackage{xcolor}
\usepackage[margin=2.5cm]{geometry}
\usepackage[hidelinks]{hyperref}
\usepackage{setspace}
\onehalfspacing

% -- Convenience macros (shared with the other chapters) -----------------
\newcommand{\dpdk}{\textsc{sd-core dpdk}}
\newcommand{\usr}{\textsc{usr-upf}}
\newcommand{\upf}{\textsc{upf}}
\newcommand{\code}[1]{\texttt{#1}}
\newcommand{\figpath}[1]{figures/#1}
\newcommand{\mappo}{\textsc{mappo}}
\newcommand{\ippo}{\textsc{ippo}}
\newcommand{\ppo}{\textsc{ppo}}

\newcommand{\drawioplaceholder}[1]{%
  \fbox{\parbox{0.94\linewidth}{\centering\vspace{2pt}
    \textbf{[\,DRAW.IO PLACEHOLDER\,]}\\[4pt]
    \footnotesize\raggedright #1\vspace{2pt}}}}

\begin{document}
% ============================================================

\chapter{Learning Controllers for Energy-Aware UPF Orchestration}
\label{chap:controllers}

% ------------------------------------------------------------------------
\section{Introduction}
\label{sec:ctrl_intro}

The preceding chapters built an \emph{observability stack}. The profiling
campaign (Chapter~\ref{chap:upf_profiling}) made per-network-function power and
Quality-of-Service (QoS) compliance observable---quantities a production 5G core
cannot instrument, because hardware energy counters are per-package, not
per-\upf{}. The forecaster (Chapter~\ref{chap:forecaster}) made the
\emph{near future} observable, projecting per-cluster load one step ahead of the
present. The digital twin (Chapter~\ref{chap:digital_twin}) made the
\emph{counterfactual} observable: what power and QoS \emph{would} result under a
\upf{} realisation that is not currently running.

This chapter asks what that augmented observability is \emph{for}, and whether it
can be trusted. The vehicle is a concrete control problem---selecting, at each
15-minute interval and at each Mobile Edge Computing (MEC) site, between a
high-throughput \dpdk{} data plane and a lower-power \usr{} data plane---and the
central claim is methodological:

\begin{quote}
Energy saving is not the contribution; it is the \emph{instrument}. A controller
trained purely on inferred signals---attributed power, forecast load,
counterfactual QoS---can only beat classical baselines on a held-out period if
those inferred signals are faithful. That it does is the validation of the
observability stack. Control performance is the measurement, not the end.
\end{quote}

The chapter follows the historical development of the controllers, because that
development \emph{is} the argument. Section~\ref{sec:ctrl_predecessor} describes
the single-site predecessor system, which acted on raw telemetry and one offline
forecast and motivated everything that followed. Section~\ref{sec:ctrl_baselines}
formalises the rule-based baselines any learning controller must beat.
Section~\ref{sec:ctrl_singlesite} introduces the single-site reinforcement
learning (RL) formulation and its result. Section~\ref{sec:ctrl_lstm} reports a
deliberate negative detour---a recurrent policy---because where the temporal
observability should live is itself a finding. Section~\ref{sec:ctrl_ceiling}
explains why a single-site view has a ceiling, and Section~\ref{sec:ctrl_marl}
lifts that ceiling with cooperative multi-agent RL. Section~\ref{sec:ctrl_synth}
draws the observability argument together.

% ------------------------------------------------------------------------
\section{The single-site predecessor}
\label{sec:ctrl_predecessor}

The work began as a single, self-contained system (the \code{EnergyAwareUPF}
prototype) that combined, in one repository, a hand-fitted power model, an
offline traffic predictor, and a \ppo{} controller for one \upf{} site. It
established the problem shape and the reward ingredients that survive into the
present work, and its limitations are precisely what drove the modular
observability stack the thesis now presents.

Three of its design choices are worth stating, because the later chapters are
reactions to them. First, its power and QoS model was fitted inside the control
repository from a bespoke measurement set; there was no separation between the
\emph{evaluator} and the \emph{controller}, so the controller could not be said
to have been tested against an independent instrument. Second, its traffic
forecast was a Keras Long Short-Term Memory (LSTM) network trained offline on a
single time series, its predictions exported to a column the environment read as
static data---not a live model in the loop. Third, it was single-site by
construction: one environment, one agent, one \upf{} location.

Quantitatively the predecessor already showed that the problem is real: across
five independently trained seeds a feed-forward \ppo{} controller reduced mean
attributed energy by $16.9\%$ relative to always-on \dpdk{} provisioning, at an
equal QoS-violation rate to a hand-tuned rule-based threshold. It also showed,
through an explicit ablation, that a \emph{recurrent} \ppo{} policy did not
improve on the feed-forward one---a result this chapter revisits in a cleaner
setting (Section~\ref{sec:ctrl_lstm}).

The predecessor is therefore not presented as current work but as the
\emph{pre-augmentation baseline}: a controller acting on raw telemetry and one
frozen forecast, with an in-repository power model. Each subsequent chapter
replaces one of its improvisations with a first-class, independently validated
observability component---and this chapter shows that doing so is what makes the
controller's results trustworthy rather than merely favourable.

% ------------------------------------------------------------------------
\section{Problem formulation and rule-based baselines}
\label{sec:ctrl_baselines}

\paragraph{Decision problem.} A fleet of $K=10$ heterogeneous MEC clusters, each
serving an aggregated traffic load derived from the NetMob Lyon traces at
15-minute resolution, must at every step select per site between \dpdk{} and
\usr{}. \dpdk{} sustains a near-constant $0.82$\,W power floor but tolerates high
load; \usr{} is far cheaper at low load but has a narrow safe operating region
and violates the delay/loss budget above roughly $0.15$\,Gbps. Switching
realisations is not free: it incurs a measured activation energy and, optionally,
a cooldown penalty. The controller sees, per site, the current load, an
eight-step load history, a one-step-ahead forecast, and the recent QoS, energy,
and switching state (a $14$-dimensional observation).

\paragraph{Baselines from the twin.} The rule-based controllers are not
hand-tuned; they are \emph{derived} from the digital twin's surrogate by
sweeping load and reading off the physical operating points
(Chapter~\ref{chap:digital_twin}):

\begin{center}
\begin{tabular}{lll}
\toprule
Operating point & Value & Definition \\
\midrule
Energy break-even & $91$\,Mbps & first load where \usr{} power $\ge$ \dpdk{} power \\
QoS limit & $149$\,Mbps & last load where \usr{} is safe \\
Decision threshold $\lambda_{\mathrm{dec}}$ & $81$\,Mbps & $\min(\text{break-even},\text{QoS}) - 10$\,Mbps margin \\
Forecast MAE ($K{=}10$) & $110.7$\,Mbps & from the forecaster's held-out summary \\
\bottomrule
\end{tabular}
\end{center}

Two families follow. A \emph{threshold} controller switches to \usr{} whenever
predicted load is below $\lambda_{\mathrm{dec}}$. A \emph{hysteresis} controller
adds a band, $\lambda_{\uparrow}=81$, $\lambda_{\downarrow}=\lambda_{\uparrow}-\text{band}$,
to suppress flapping under forecast noise. Note a subtlety the twin chapter also
records: the paper-faithful automatic band is $2\times\text{MAE}=221$\,Mbps,
which is \emph{wider} than the decision point itself and collapses the controller
to always-\dpdk{}; a tuned narrow band ($20$\,Mbps) is the real classical
contender. This distinction---between the derived-but-degenerate rule and the
tuned rule---is exactly the kind of brittleness that motivates learning.

\paragraph{Reward.} All learning controllers optimise
\begin{equation}
  r_t \;=\; -\bigl(\alpha\,\mathrm{SEC}_t \;+\; L^{\mathrm{QoS}}_t \;+\; L^{\mathrm{SW}}_t \;+\; L^{\mathrm{CD}}_t\bigr),
  \label{eq:ctrl_reward}
\end{equation}
where $\mathrm{SEC}_t$ is the steady-state specific energy consumption (power per
Mbps), $L^{\mathrm{QoS}}_t=\lambda_{\mathrm{QoS}}\max(0,\tau-Q_t)$ is a graded
penalty on the twin's continuous QoS score $Q_t$, $L^{\mathrm{SW}}_t=\lambda_{\mathrm{sw}}\cdot
\text{sw\_energy\_wh}$ charges the twin's measured switching energy only when the
realisation actually changes, and $L^{\mathrm{CD}}_t$ is a soft cooldown
surcharge. Crucially the switching energy is a \emph{load-independent per-variant
constant}---activation happens at zero traffic---so the switch penalty does not
spuriously scale with the load the new \upf{} has not yet begun to serve
(Chapter~\ref{chap:digital_twin}). Every term in Equation~\eqref{eq:ctrl_reward}
is a reading from the observability stack: SEC and $Q$ from the twin's
surrogates, the forecast inside the observation from the forecaster, the
switching energy from the profiling campaign.

% ------------------------------------------------------------------------
\section{Single-site reinforcement learning}
\label{sec:ctrl_singlesite}

The first learning controller is a feed-forward \ppo{} agent on a single cluster,
trained on the historical (train) split, checkpoint-selected on the validation
split, and evaluated once on the held-out test split. On the representative
low-load cluster it reduces energy and total penalty substantially relative to
both a static policy and the derived threshold:

\begin{center}
\begin{tabular}{lrrr}
\toprule
Policy & Total reward & Energy (Wh) & Unsafe \% \\
\midrule
Single-site \ppo{} & $\mathbf{-881.1}$ & $\mathbf{176.6}$ & $0.37$ \\
Always \dpdk{} & $-1321.4$ & $206.8$ & $0.00$ \\
Threshold (derived $81$\,Mbps) & $-1348.1$ & $175.9$ & $2.48$ \\
\bottomrule
\end{tabular}
\end{center}

The \ppo{} controller improves total reward by $33.3\%$ and energy by $14.6\%$
over always-\dpdk{}, and beats the derived threshold on every axis at once---it
is both safer (0.37\% vs 2.48\% unsafe) and cheaper. The reason is observability:
the threshold acts on the noisy one-step forecast alone, so it flips on forecast
jitter, whereas the learner also sees the eight-step history, which acts as an
implicit smoother, and the twin's QoS and energy signals, which let it weigh the
cost of a wrong switch. The single-site result is the first evidence that the
inferred signals are good enough to act on.

\begin{figure}[t]
  \centering
  \includegraphics[width=0.82\linewidth]{\figpath{fig_p2_singlesite_reward.png}}
  \caption{Single-site \ppo{} against static and threshold baselines on the
  held-out test period for a low-load cluster. The learner dominates on both
  energy and total penalty (cluster~0, \code{v0.4} twin).}
  \label{fig:ctrl_singlesite}
\end{figure}

% ------------------------------------------------------------------------
\section{A detour: where should temporal observability live?}
\label{sec:ctrl_lstm}

The observation vector encodes time explicitly, through an eight-step load
history. An alternative is to let the policy network carry temporal state
implicitly, through recurrence---an LSTM policy that reads only the current
observation and maintains its own hidden memory. Both the predecessor system and
a controlled ablation in this work tested this. The finding is consistent and,
for a thesis about observability, informative: the recurrent policy did
\emph{not} improve on the feed-forward policy with an explicit history window,
and cost more to train.

The interpretation matters more than the result. It says the temporal
observability the controller needs is better supplied \emph{as an explicit,
inspectable feature of the observed state} than manufactured inside an opaque
recurrent policy. This is the same design principle the whole stack embodies:
make the decision-relevant signal an explicit, validated observable rather than
an implicit internal quantity. The explicit history window is retained in every
subsequent controller, and the recurrent variant is not pursued further.

% ------------------------------------------------------------------------
\section{The single-site ceiling}
\label{sec:ctrl_ceiling}

A per-site controller, whether rule-based or learned, sees only its own site.
That is a real limitation once the deployment is a fleet of heterogeneous
clusters. Some clusters carry dense, persistent demand and generate abundant
experience around the switching boundary; others sit at low load for long
periods and rarely enter the unsafe region at all. An independent learner at a
low-load, data-poor site simply does not observe enough boundary-crossing events
to identify \emph{when} \usr{} stops being safe. Its own history is not
informative enough.

Two multi-site designs bracket the problem. A single \emph{centralised} \ppo{},
observing the concatenated $140$-dimensional fleet state and emitting a joint
$10$-dimensional action, has the information in principle but must learn one
monolithic policy over a joint space; at a fixed training budget it is fragile.
In this work the centralised controller collapsed---on one low-load cluster
(\code{c6} in the illustrated seed) it degenerated to near-always-\usr{}, driving
that cluster's QoS-violation rate well above every other controller's, and its
across-seed reward variance was an order of magnitude larger than any other
controller (Figure~\ref{fig:ctrl_collapse}). The opposite design, an
\emph{independent} \ppo{} per site (an \ippo{} ensemble), is stable and
embarrassingly parallel to train, but by construction cannot share what one site
learns with another---so the data-poor clusters remain data-poor.

\begin{figure}[t]
  \centering
  \includegraphics[width=0.9\linewidth]{\figpath{fig_p3_centralised_collapse.png}}
  \caption{Per-cluster, per-step action of the centralised \ppo{} on the test
  period (green = \dpdk{}, red = \usr{}). The solid red band (cluster \code{c6}
  in this seed) is a low-load cluster collapsed to near-always-\usr{}: the
  monolithic joint policy could not factor the per-site substructure at this
  training budget.}
  \label{fig:ctrl_collapse}
\end{figure}

The single-site ceiling is thus not a tuning problem but a \emph{structural} one:
the missing capability is the ability to make one site's operating experience
observable to another. That is a spatial extension of the observability
envelope, and it is what cooperative multi-agent learning supplies.

% ------------------------------------------------------------------------
\section{Cooperative multi-agent control (MAPPO)}
\label{sec:ctrl_marl}

The resolution is centralised training with decentralised execution (CTDE),
realised as multi-agent \ppo{} (\mappo{}). A single \emph{shared} actor network
maps one site's $14$-dimensional observation to that site's action; the same
parameters serve all $K=10$ agents, so each gradient step aggregates experience
across the fleet. A \emph{centralised} critic sees the full $140$-dimensional
joint state during training and produces a value estimate per agent, so credit
assignment respects the shared reward structure. Execution remains fully
decentralised: at run time each site acts on its own observation alone.

\begin{figure}[t]
  \centering
  \includegraphics[width=0.86\linewidth]{\figpath{fig_p7_reward_errorbars.png}}
  \caption{Test-period reward, mean over eight training seeds, with standard
  deviation, under the \code{v0.4} twin. \mappo{} is the best controller and the
  most stable trained one; values match Table~\ref{tab:ctrl_headline}.}
  \label{fig:ctrl_errorbars}
\end{figure}

\paragraph{Headline result.} Across eight independent training seeds, evaluated
once each on the held-out test period under the corrected (\code{v0.4}) twin:

\begin{table}[t]
\centering
\caption{Test-period comparison across all controllers (mean $\pm$ std over
seeds where trained). \mappo{} is best on reward and QoS at equal energy.}
\label{tab:ctrl_headline}
\begin{tabular}{lrrrr}
\toprule
Controller & Reward & Energy (Wh) & Unsafe \% & Seeds \\
\midrule
\mappo{} (CTDE) & $\mathbf{-5647 \pm 86}$ & $1967$ & $\mathbf{0.53}$ & $8$ \\
\ippo{} ensemble & $-5789 \pm 83$ & $1974$ & $0.57$ & $8$ \\
Hysteresis (band $20$, cd $1$) & $-6544$ & $1971$ & $0.94$ & det. \\
Always \dpdk{} & $-7201$ & $2071$ & $0.38$ & det. \\
Threshold (derived $81$\,Mbps) & $-7623$ & $1970$ & $1.42$ & det. \\
Centralised \ppo{} & $-14459 \pm 4509$ & $2130$ & $3.73$ & $4$ \\
\bottomrule
\end{tabular}
\end{table}

\mappo{} beats every baseline, including the best hand-tuned classical
controller, at the lowest QoS-violation rate of any trained controller and at
energy indistinguishable from the ensemble. The margin over the \ippo{} ensemble
($\approx 141$ reward units) exceeds the combined seed standard deviation, and a
paired bootstrap over the eight seeds confirms the ordering is significant:
mean $\Delta = 141$, $95\%$ CI $[23, 237]$ (excludes zero), $p = 0.011$
(\mappo{}${}>{}$\ippo{}), computed under the \code{v0.4} twin.

\paragraph{Where the win comes from.} The advantage is not uniform---it
concentrates exactly where the single-site ceiling predicted it would.
Figure~\ref{fig:ctrl_bridge} places the independent ensemble and \mappo{}
side by side per cluster. On high-load clusters (c3, c9) every controller
correctly stays on \dpdk{} and the bars coincide. The gains appear on the
mid- and low-load, data-poor clusters (c0, c4, c5, c7), where the independent
learner could not accumulate enough boundary experience to justify \usr{}, but
\mappo{}'s shared actor---trained simultaneously on the data-rich clusters---
transferred the ``\usr{} is safe at low predicted load'' rule in and applied it.
This transfer is the spatial observability extension made concrete: knowledge
observed at one site becomes actionable at another.

\begin{figure}[t]
  \centering
  \includegraphics[width=\linewidth]{\figpath{fig_bridge_ippo_vs_mappo.pdf}}
  \caption{Per-cluster test reward, independent single-site learners (\ippo{}
  ensemble) versus cooperative \mappo{}, mean over eight seeds under the
  \code{v0.4} twin. Shaded clusters are where \mappo{} exceeds the ensemble by
  more than $15$ reward units---the data-poor sites that inherit the operating
  rule from the fleet.}
  \label{fig:ctrl_bridge}
\end{figure}

\begin{figure}[t]
  \centering
  \includegraphics[width=0.9\linewidth]{\figpath{fig_mappo_action_heatmap.png}}
  \caption{\mappo{} per-site, per-step actions on the test period. The low-load
  clusters show the intelligent ``\usr{} at low predicted load'' pattern that
  the collapsed centralised policy (Figure~\ref{fig:ctrl_collapse}) could not
  find, now learned in the joint setting.}
  \label{fig:ctrl_heatmap}
\end{figure}

\paragraph{Stability and safety.} \mappo{} is also the most seed-stable trained
controller, and the safety--reward trade-off separates the architectures
cleanly: the trained cooperative and independent controllers cluster in the
high-reward, low-unsafe corner, while the centralised policy scatters along the
safety axis---visual evidence that the wrong coordination structure carries a
real safety cost, not merely a reward cost (Figure~\ref{fig:ctrl_safety}).

\begin{figure}[t]
  \centering
  \includegraphics[width=0.72\linewidth]{\figpath{fig_p7_safety_vs_reward.png}}
  \caption{Safety versus reward on the test period. Top-left is ideal.
  Cooperative and independent learners (eight seeds each) cluster there;
  the centralised policy spreads along the unsafe axis.}
  \label{fig:ctrl_safety}
\end{figure}

% ------------------------------------------------------------------------
\section{Synthesis: control as validation of the observability stack}
\label{sec:ctrl_synth}

The controllers form a ladder, and each rung consumes a wider slice of the
observability envelope than the one below it. The derived threshold acts on a
single inferred signal, the one-step forecast, and is brittle precisely because
that is all it sees. The single-site learner adds the twin's counterfactual QoS
and energy signals and an explicit temporal history, and beats the threshold on
every axis. The cooperative fleet controller adds spatial observability---one
site's experience made visible to another---and beats the best classical
controller by a margin larger than seed noise, at the lowest QoS-violation rate
of any learner.

The point of the ladder is not that more reinforcement learning is better. It is
that each measurable improvement in control tracks a specific extension of what
the system can observe, and that the improvements only materialise because the
inferred signals are faithful. A controller trained on wrong attributed power, a
wrong forecast, or a wrong counterfactual could not have beaten a controller
tuned on the physical baselines over a held-out period---and yet the learned
controllers do, consistently and significantly. That is the sense in which the
control results \emph{validate} the observability stack rather than merely
exploit it.

Two caveats keep the claim honest. First, the load calibration relating the
dimensionless NetMob traffic to physical throughput is a declared modelling
assumption, not a measured constant; the operating regime the whole comparison
lives in follows from it (Chapter~\ref{chap:digital_twin}). Second, the
evaluation is a faithful offline replay against a measurement-grounded twin, not
a live testbed deployment; the intended next step is closed-loop operation with
the controller as an analytics-driven decision function alongside standardised
5G analytics functions. Neither caveat undermines the central finding: augmented
observability, built from profiling, forecasting, and counterfactual simulation,
is accurate enough that a controller can act on it and outperform control that
waits for the physical ground truth.

\end{document}
